\section{Limits of XOR-frame recursions}\label{sec:frames}

The only technique claimed to beat $O(L\log n)$ for transposition on tapes is an unrefereed
construction of OpenAI~\cite{OpenAI2026mult}, which claims $O(L(\log n)^{1-\kappa})$ for a tiny
$\kappa>0$. Its core is an address-interchange primitive that mixes data by XOR in sheared address
frames, applied recursively. This section abstracts the recursion and shows that it cannot give linear
time: every gadget it could use makes at least as many frame changes as the trivial one, so its cost is at
least $W$. The results are about an
abstract model, so they rule out a class of constructions, not linear-time transposition.

\subsection{The framework}

Fix $m\ge2$. Two address chunks $H$ and $D$ are to be interchanged; write each as a vector of $m$
fields. For a rational $m\times m$ matrix $M$ the \emph{frame} $\Phi_M$ is the address map
$(H,D)\mapsto(H+MD,D)$; frames compose additively, $\Phi_A\Phi_B=\Phi_{A+B}$. Changing the frame of a
stream from $A$ to $B$ costs $\rank(B-A)$ recursive interchanges of narrower chunks.

\begin{definition}[Gadget]\label{def:gadget}
\formal{Frames/Gadget.lean}{Transpose.ShearFrame.Gadget, Transpose.ShearFrame.Gadget.cost}
A \emph{gadget} has $W$ wires; wire $w$ starts in a source frame $S_w$, passes through the frames of its
gates in time order, and ends in a sink frame $T_w$. Gates are XOR operations between wires in a common
frame. The scalar network permutes the wires by a permutation $\rho$ (G1), and
$T_{\rho(w)}=S_w+I$ (G2): every wire ends sheared by the identity. The \emph{cost} $s$ of the gadget is
the sum over wires of the ranks of consecutive frame differences along the wire, from source to sink.
Formally, \lean{Transpose.ShearFrame.Gadget} records $\rho$, the source frames and the frame paths subject to
(G2); the gates and (G1) enter only through the run model of \cref{def:frame-runs}.
\end{definition}

The recursion built on a gadget distributes a stream over the $W$ wires and runs the gadget, each frame
change being a recursive interchange of narrower chunks. Its cost per unit of data satisfies
$F_k\le(s/W)F_{k-1}+O(1)$ at width $m^k$, and every implementation in the template also satisfies
$F_k\ge(s/W)F_{k-1}+c_0$. A gadget with $s<W$ would therefore give linear-time transposition, a gadget
with $s/W=m^\tau$, $0<\tau<1$, would give $O(L(\log n)^\tau)$, and $s\ge W$ excludes linear time within the
template; $s=W$ gives $\Theta(L\log\log n)$. (This reading of the template is informal: it rests on the
preprint's construction~\cite{OpenAI2026mult}, which we use only to motivate the model.)

\subsection{Cost bounds}

\begin{theorem}[Cycle bound and projection bound]\label{thm:cycle}
\formal{Frames/Gadget.lean}{Transpose.ShearFrame.Gadget.invariant_cost_ge, Transpose.ShearFrame.Gadget.cost_ge_of_invariant_family, Transpose.ShearFrame.Gadget.cost_ge_m}
\formal{Frames/GadgetProjection.lean}{Transpose.ShearFrame.PGadget.two_mul_cost_ge}
Every nonempty $\rho$-invariant set of wires has total cost at least $m$, so $s\ge jm$ for $j$ disjoint
such sets, and $s\ge m$ if $W\ge1$. If every wire passes through at least one gate and every gate frame is
idempotent (a projection, as in the preprint), then $2s\ge Wm$.
\end{theorem}

\begin{outline}
Rank is subadditive, so the cost of a wire is at least $\rank(T_w-S_w)$, and over a cycle $C$ of $\rho$,
$\sum_{w\in C}(T_w-S_w)=\abs{C}I$ has rank $m$. For idempotent frames, $\rank(P_V-P_U)\ge\abs{\dim V-\dim
U}$ and $\rank(P_A+P_B)\ge\max(\dim A,\dim B)$; pairing the source term of $w$ with the sink term of
$\rho(w)$ and optimizing gives $2s\ge Wm$.
\end{outline}

The general lower bound uses a run model in which frames are arbitrary and only frame changes are counted.

\begin{definition}[Runs]\label{def:frame-runs}
\formal{Frames/GadgetRuns.lean}{Transpose.ShearFrame.Cfg, Transpose.ShearFrame.Event, Transpose.ShearFrame.moves}
A configuration assigns to each of $W$ wires a frame (from any set $F$) and a vector in $\F_2^W$. An
event is either a gate $i\leftarrow i+j$, allowed only for $i\ne j$ at a common frame, or a move of one
wire to a new frame. A run starts with wire $u$ at frame $S_u$ holding the unit vector $e_u$.
\end{definition}

\begin{theorem}[At least $W$ frame changes]\label{thm:frame-change}
\formal{Frames/GadgetRuns.lean}{Transpose.ShearFrame.moves_ge_W, Transpose.ShearFrame.cost_ge_W, Transpose.ShearFrame.moves_ge_W_matrix}
Let $\mathrm{shift}:F\to F$ have no fixed point. Every valid run that ends with wire $\rho(u)$ at frame
$\mathrm{shift}(S_u)$ holding $e_u$, for every $u$, contains at least $W$ moves. Hence every gadget has cost $s\ge W$ (for matrix frames, $\mathrm{shift}(S)=S+I$): reading its gates as
gate events and its frame changes as moves, (G1) and (G2) say that the run ends with wire $\rho(u)$ at
$S_u+I$ holding $e_u$, and a genuine frame change costs the rank of a nonzero matrix, at least $1$ (this
translation is not formalized; the structure of \cref{def:gadget}, which omits the gates and (G1), only
satisfies $s\ge m$). So the framework cannot give linear time.
\end{theorem}

\begin{outline}
Let $\sigma_f$ be the span of the vectors on the wires at frame $f$, and $N_f$ the span of the $e_u$ with
$\mathrm{shift}(S_u)=f$. The potential $\sum_f\dim(\sigma_f\cap N_f)$ is unchanged by gates (they act
within one frame by transvections), grows by at most one per move, is $0$ initially (no unit vector starts
at its target, since $\mathrm{shift}$ has no fixed point), and is $W$ at the end.
\end{outline}

\begin{theorem}[Mixing across frames does not help]\label{thm:group-algebra}
\formal{Frames/GroupAlgebraRuns.lean}{Transpose.ShearFrame.GroupAlgebra.changes_ge_W, Transpose.ShearFrame.GroupAlgebra.cost_ge_W}
In the group-algebra model, where each wire holds a formal sum of sheared items (a function from an
additive group of frames to $\F_2^W$), gates $i\leftarrow i+j$ are allowed between any wires, and a change
shifts the content of one wire, every run from the configuration in which wire $u$ holds exactly item $u$ at frame $S_u$ to one in which
wire $\rho(u)$ holds item $u$ at frame $S_u+I$ and nothing else ($I\ne0$) makes at least $W$ changes.
\end{theorem}

\begin{outline}
The span of the vectors $w\mapsto(\text{content of }w\text{ at frame }I+S_u)_u$ has dimension $0$ at the
start and $W$ at the end; gates do not increase it and each change increases it by at most one.
\end{outline}

\subsection{Coarsening and integer frames}

\begin{theorem}[Coarsening]\label{thm:coarsening}
\formal{Frames/GadgetCoarse.lean}{Transpose.ShearFrame.vmoves_ge_W, Transpose.ShearFrame.visible_ge_W_of_coarse, Transpose.ShearFrame.det_ne_zero_of_moves_le, Transpose.ShearFrame.rank_sum_of_moves_le}
For every map $\pi:F\to F'$ with $\pi(\mathrm{tgt}(u))\ne\pi(S_u)$ for all $u$, every valid run as in
\cref{thm:frame-change} (with targets $\mathrm{tgt}(u)$) makes at least $W$ moves that change the
$\pi$-class of the moved wire. For rational matrix frames with targets $S_u+I$, a run with at most $W$
moves changes every frame by a matrix of nonzero determinant, and the total rank of its jumps is at least $Wm$ (for a gadget whose wires each change frame once this
means cost at least the baseline $Wm$, through the same unformalized translation as in
\cref{thm:frame-change}).
\end{theorem}

\begin{outline}
The potential of \cref{thm:frame-change} needs only a target frame per item and survives coarsening: a move
that does not change the $\pi$-class of the wire does not change the coarsened run. For single jumps: if a jump $J$ of a run with at most $W$ moves were singular, choose $\ell\ne0$ with
$\ell^{\T}J=0$; the coarsening $M\mapsto\ell^{\T}M$ separates every $S_u$ from $S_u+I$ (they differ by
$\ell^{\T}\ne0$) but hides $J$, so fewer than $W$ moves would be visible. So every jump has rank $m$, and
there are at least $W$ jumps by \cref{thm:frame-change}.
\end{outline}

\begin{theorem}[Mod-$2$ bounds for integer frames]\label{thm:mod2}
\formal{Frames/GadgetMod2.lean}{Transpose.ShearFrame.two_mul_le_rank_sum}
For integer $2\times2$ frames, the total rank of the frame jumps of every valid run that ends with wire
$\rho(u)$ at frame $S_u+I$ holding $e_u$ (as in \cref{thm:frame-change}) is at least $2W$:
gadgets with $2\times2$ integer frames never beat the baseline.
\end{theorem}

\begin{theorem}[Mod-$2$ bounds, every $m$]\label{thm:mod2m}
\formal{Frames/GadgetOneSided.lean}{Transpose.ShearFrame.one_sided_bound}
\formal{Frames/GadgetTwoSided.lean}{Transpose.ShearFrame.two_sided_bound}
For integer $m\times m$ frames, the total rank $s$ (over $\Q$) of the jumps of every valid run that ends with
wire $\rho(u)$ at frame $S_u+I$ holding $e_u$ satisfies
$2(2^m-1)W\le2^ms$ and $2(2^m-1)W\le(2^{m-1}+1)s$. In particular $s\ge2.8W$ for $m=3$ and
$s\ge(4-o(1))W$ as $m\to\infty$.
\end{theorem}

\begin{outline}[Outline of the formal proofs of \cref{thm:mod2,thm:mod2m}]
Average \cref{thm:coarsening} over a family of coarsenings: $M\mapsto\ell^{\T}(M\bmod2)$ for nonzero
$\ell\in\F_2^m$ (one-sided), or $M\mapsto\ell^{\T}(M\bmod2)v$ for pairs with $\ell^{\T}v=1$ (two-sided). Each
coarsening separates every source from its target, so it sees at least $W$ jumps. A jump of rank $r$ over
$\F_2$ is seen by $2^m-2^{m-r}$ of the one-sided coarsenings, and by a bounded number of the two-sided ones,
computed from the left kernel and the fixed vectors of the jump; and the rank over $\F_2$ is at most the
rank over $\Q$ (\leanfile{Frames/RankMod.lean}). For $m=2$ a finite check over $\F_2^{2\times2}$ gives the
constant.
\end{outline}

\subsection{Tight runs}

A valid run that takes every item $u$ from its source frame $S_u$ to a target frame $\mathrm{tgt}(u)\ne S_u$
with at most $W$ frame-changing moves is \emph{tight}; by \cref{thm:coarsening} it makes exactly $W$ such moves.
Tight runs are rigid.

\begin{theorem}[Structure of tight runs]\label{thm:tight}
\formal{Frames/GadgetCargo.lean}{Transpose.ShearFrame.tight_potentials, Transpose.ShearFrame.tight_sigma_split, Transpose.ShearFrame.cargo_rule, Transpose.ShearFrame.cargo_rule_bool}
In a tight run, after any prefix with $t$ frame-changing moves the ``arrived'' potential is $t$ and the ``not yet
departed'' potential is $W-t$, and at every frame the span of the vectors there splits into its arrived and its
not-departed parts. If moreover each item has only two frames (every frame other than $S_u$ is
$\mathrm{tgt}(u)$, as after a two-class coarsening), then every frame-changing move carries a vector lying in the
span of the items whose target is its destination and in the span of the items born at its origin (the cargo
rule). (The formal statements assume finitely many frames; a run uses only finitely many, but this transfer
is not formalized.)
\end{theorem}

\begin{theorem}[Item itineraries]\label{thm:itineraries}
\formal{Frames/GadgetItinerary.lean}{Transpose.ShearFrame.chain_exists_visible, Transpose.ShearFrame.itinerary_condition, Transpose.ShearFrame.chain_itinerary, Transpose.ShearFrame.itinerary_condition_matrix, Transpose.ShearFrame.chain_itinerary_matrix}
In every valid run as in \cref{thm:coarsening}, every item has a chain of frame-changing hops (moves whose cargo
contains the item) from its source to its target. Let $\pi$ be a two-class coarsening with
$\pi(\mathrm{tgt}(u))\ne\pi(S_u)$ for all $u$, and suppose the run makes exactly $W$ moves that change the
$\pi$-class. Then every such move carrying an item starts on the item's birth side ($\pi$ of its origin equals
$\pi(S_u)$), and every item has a chain from its source to its target all of whose $\pi$-visible hops are
admissible in this sense. For frames in $\F_2^{m\times m}$ with targets $S_u+I$ and a coarsening by a functional
$\Lambda$ with $\operatorname{tr}\Lambda=1$ (the class of $M$ is $\langle\Lambda,M\rangle$), admissibility of a hop of $s$
from $f_0$ with jump $J$ and $\langle\Lambda,J\rangle=1$ reads $\langle\Lambda,S_s+f_0\rangle=0$.
\end{theorem}

\begin{theorem}[Every carrying move lies on an admissible chain]\label{thm:carrying}
\formal{Frames/GadgetItineraryC.lean}{Transpose.ShearFrame.carrying_move_on_chain, Transpose.ShearFrame.carrying_move_on_admissible_chain, Transpose.ShearFrame.visible_move_on_admissible_chain}
Every frame-changing move whose cargo contains an item lies on a chain of visible hops of that item from its
source to its target; under tightness, every visible move lies on an admissible chain of some item.
\end{theorem}

\begin{theorem}[Mixing lemma]\label{thm:mixing}
\formal{Frames/GadgetMixing.lean}{Transpose.ShearFrame.mixed_move_forces_extra_hop, Transpose.ShearFrame.mixing_lemma, Transpose.ShearFrame.mixing_lemma_matrix}
Weight frames by a function $d$ with $d(f,f)=0$, $d(f,g)=0\Rightarrow f=g$ and the triangle inequality (for
matrices, $d(A,B)=\rank(B-A)$), and let $d(S_u,\mathrm{tgt}(u))\ge m$ for every item $u$. Suppose a
valid run as in \cref{thm:frame-change} contains a frame-changing move whose cargo contains two distinct
items, and every other frame-changing move carries a single item. Then some item of that cargo has
$d$-weighted hop count at least $m+1$: a mixed move forces an extra hop.
\end{theorem}

\begin{outline}[Outline of the formal proofs of \cref{thm:tight,thm:itineraries,thm:carrying,thm:mixing}]
The vectors of a configuration reached by a valid run are linearly independent, since gates are
transvections; so the spans at distinct frames are independent and their dimensions add up to $W$. The
arrived and not-departed potentials each change by at most one per visible move and sum to at most $W$, so in
a tight run both move in lockstep, which forces the split and the cargo rule. Chains are found by following an
item backwards (and forwards) through the moves whose cargo contains it, using that an item is present at a
frame only if some move carried it there. Admissibility (\cref{thm:itineraries}) is the cargo rule of the coarsened two-class run. For the mixing
lemma, suppose every item of the mixed cargo $x$, moved into frame $g$, has hop count at most $m$. An item
present at a frame after (before) a segment was (will be) present at a frame within $d$-distance of it
bounded by the segment's hop count; with the triangle inequality, no item of the support of $x$ is present at
$g$ before the mixed move, and no later frame-changing move into $g$ carries one. As the later moves out of
$g$ carry unit vectors, every wire at $g$ stays constant on the support of $x$ while some wire at $g$ keeps
coordinate $s_0$; at the end that wire holds $e_{s_0}$, which is not constant on a support of two items.
\end{outline}

Together these say that, within the XOR-frame template, every gadget costs at least $W$ (so linear time is
out of reach), gadgets with integer frames cost substantially more than $W$, and gadgets that attain the bound (tight runs) are rigid. Whether some gadget has $s/W$ substantially below $m$, which would give an improvement
over $O(L\log n)$ much larger than the claimed $\kappa$, is open (\cref{sec:open}).
